\subsection{Every adjacent norm in the pure seventeen tower}

\begin{statement}{E.65}{Proved}{Pure norm images, an explicit wild character, and the certified first unit norm}
Let
\[
 F_d=\mathbb Q(\alpha_d),\qquad \alpha_0=17,\qquad
 \alpha_{d+1}^{,3}=\alpha_d,\qquad
 M_d=F_d(\zeta),\quad \zeta=\zeta_3,
 \tag{E65.1}
\]
where the $\alpha_d$ are positive real roots. Complex conjugation $c$
fixes $F_d$ and sends $\zeta$ to $\zeta^{-1}$.
Write $\widehat{F_d^\times}_3=\varprojlim_m
F_d^\times/(F_d^\times)^{3^m}$, with the multiplication-completion
topology. The following formulas evaluate the image and every fiber of
the actual adjacent norm
\[
 N_d^{\rm p}:\widehat{F_{d+1}^\times}_3
                    \longrightarrow\widehat{F_d^\times}_3.
 \tag{E65.2}
\]

For a prime $v\nmid3\cdot17$ of $F_d$, let $q_v=\mathrm Nv$.
When $q_v\equiv1\pmod3$, choose a primitive third root
$\zeta_v\in\mathbb F_{q_v}$, and define
\[
 \overline{\alpha_d}^{\,(q_v-1)/3}
             =\zeta_v^{\gamma_{d,v}},\qquad
 \mathcal I_d^{\rm p}=
 \{v\nmid3\cdot17:q_v\equiv1\pmod3,
                         \ \gamma_{d,v}\ne0\}.
 \tag{E65.3}
\]
The choice of $\zeta_v$ specifies one of the two primes of $M_d$ over
$v$. Changing that choice multiplies this one coordinate by $-1$.
For $v\in\mathcal I_d^{\rm p}$ set
$\delta_{d,v}(z)=\gamma_{d,v}v(z)\pmod3$.

For $d\ge1$ there is one further, explicitly computable character
$\delta_{d,3}:F_d^\times\to\mathbb F_3$.
Its finite matrix definition below uses the completion of $F_d$ at
its unique prime of ramification index $2\cdot3^{d-1}$ over three.
There is no such coordinate for $d=0$.
With that convention there is an exact sequence of topological groups
\[
 \widehat{F_{d+1}^\times}_3
 \xrightarrow{\ N_d^{\rm p}\ }
 \widehat{F_d^\times}_3
 \xrightarrow{\ \Theta_d^{\rm p}\ }
 \begin{cases}
   \displaystyle\bigoplus_{v\in\mathcal I_0^{\rm p}}\mathbb F_3,
                                                        &d=0,\\[2mm]
   \displaystyle\mathbb F_3\oplus
                \bigoplus_{v\in\mathcal I_d^{\rm p}}\mathbb F_3,
                                                        &d\ge1,
 \end{cases}
 \longrightarrow0,
 \tag{E65.4}
\]
where $\Theta_d^{\rm p}(z)$ consists of $\delta_{d,3}(z)$, when
present, and the $\delta_{d,v}(z)$. The right-hand side is discrete.
Every coordinate kills cubes, so the map on the completion is
evaluated on any actual representative modulo cubes. Its prime
coordinates have finite support. Thus (E65.4) is a sufficient as well
as necessary global norm test, and its norm image is open.

Here is the finite definition of the wild character, including all
arithmetic choices. There is a unique $\rho\in\mathbb Q_3$ with
$\rho^3=17$; it satisfies $\rho\equiv23\pmod{27}$.
For $d\ge1$ put
\[
 m=3^{d-1},\qquad e=2m,\qquad n=2e+1,
 \qquad
 Q_d(Z)=(Z-1)^{2m}+\rho(Z-1)^m+\rho^2.
 \tag{E65.5}
\]
At the specified prime, $\pi=\alpha_d+1$ is a uniformizer,
$Q_d(\pi)=0$, and its integer ring is $\mathbb Z_3[\pi]$.
Its quotient modulo $\pi^n$ is exactly the finite ring
\[
 R_d=(\mathbb Z/27)[Z]/(Q_d(Z)\bmod27,\ Z^n).
 \tag{E65.6}
\]
Let $g_i=1+\pi^i$, $1\le i<n$. Every unit of $R_d$ has a unique
expression
$\epsilon\prod_{i=1}^{n-1}g_i^{a_i}$ with
$\epsilon\in\{1,-1\}$ and $a_i\in\{0,1,2\}$.
Write $\mathbf a(u)=(a_1,\ldots,a_{n-1})$.
In the finite free cubic algebra
\[
 R'_d=R_d[\Pi]/(\Pi^3-3\Pi^2+3\Pi-\pi),
 \tag{E65.7}
\]
compute norms as determinants of multiplication in the basis
$(1,\Pi,\Pi^2)$.
Form the matrix over $\mathbb F_3$ whose rows are
\[
 \mathbf a(g_i^3)\quad(1\le i<n),\qquad
 \mathbf a\bigl(\det(1+\Pi^j\mid R'_d)\bigr)
                                  \quad(1\le j<3n).
 \tag{E65.8}
\]
Its nullspace has dimension one. Let $\mathbf w_d$ be its unique
generator whose first nonzero coordinate is $1$. The specified wild
character is
\[
 \delta_{d,3}(z)=
    \mathbf w_d\cdot
     \mathbf a\bigl(z\pi^{-v_3(z)}\bmod\pi^n\bigr)
       \quad\hbox{in }\mathbb F_3,
 \tag{E65.9}
\]
where $v_3(\pi)=1$ at this one prime.
The notation $\det(1+\Pi^j\mid R'_d)$ in (E65.8) means the
determinant of multiplication by that element, not a determinant of
a scalar matrix. Formulas (E65.5)--(E65.9) give a finite algorithm for
every $d$, using only binomial coefficients, integral lattice reduction
and $3$-by-$3$ determinants.

Each nonempty fiber of (E65.2) is a torsor under the group
\[
 \mathcal R_3=
 \left\{(a_i)_{i\ge1}\in\prod_{i\ge1}\mathbb Z_3:
   \text{for every }r,\ a_i\in3^r\mathbb Z_3
                    \text{ for all but finitely many }i\right\},
 \tag{E65.10}
\]
with its multiplication-completion topology. Explicitly, choose an
actual free basis $u_i$ of the ordinary norm-one group and a lift $y_0$
of the target. The isomorphism is
$ (a_i)\mapsto y_0\prod_i u_i^{a_i}$.
In particular the ordinary norm-one group has no torsion and is free
abelian of countably infinite rank; (E65.10) retains its infinite
discrete quotients.

The first wild matrix and its certified unit witness are
\[
 \mathbf w_1=(1,2,0,0),\qquad
 \varepsilon=7\alpha_1-18,
 \qquad\mathbf a(\varepsilon)=(2,0,2,2),
 \qquad\delta_{1,3}(\varepsilon)=2.
 \tag{E65.11}
\]
Thus $\varepsilon$ is not in
$N_{F_2/F_1}(F_2^\times)$, nor in the image of its completed norm.
Both $F_1$ and $F_2$ have class number one.
The following actual fundamental units of $F_2$, written with
$x=\alpha_2$, have norm row $(0,3,0,0)$ in the fundamental-unit
basis $\varepsilon$ of $F_1$:
\[
\begin{aligned}
 u_1={}&(20+18x-5x^2-8x^3+5x^5+2x^6-3x^7-2x^8)/3,\\
 u_2={}&-18+7x^3,\\
 u_3={}&-1-4x+6x^2-6x^4+2x^5-2x^7+2x^8,\\
 u_4={}&\tfrac13\bigl(-235+148x+182x^2-50x^3-97x^4\\
       &\hspace{3em}{}-5x^5+53x^6+16x^7-25x^8\bigr),\\
 (Nu_1,Nu_2,Nu_3,Nu_4)={}&(1,\varepsilon^3,-1,1).
\end{aligned}
 \tag{E65.12}
\]
Consequently the ordinary unit norm image is
$\langle-1,\varepsilon^3\rangle$, of index three, and its kernel is
free on $u_1,-u_3,u_4$. This specifies the preimages and signs in
addition to the index.

Under E.46, the actual positional transition of E.62 is the product
of (E65.2) and (E64.2):
\[
 \widehat{F_{d+1}^\times}_3\times
       \widehat{L_{d+1}^\times}_3
  \xrightarrow{\ (N_d^{\rm p},N_d)\ }
 \widehat{F_d^\times}_3\times\widehat{L_d^\times}_3.
 \tag{E65.13}
\]
Its total norm to $\mathbb Q$ is preserved. Therefore the
$J=(3,T-1)\mathbb Z_3[[T-1]]$ condition from E.46 is preserved by
every preimage supplied by these two norm tests.
The two tests and their fiber coordinates evaluate every adjacent
positional arrow. The compatible inverse limit of these arrows is
retained as the further global section problem; an independent
selection of preimages at each finite level is not asserted to be a
compatible selection.
The manuscript anchor is \S51.16, scans172--178,
folios733v--736v, equations(371)--(379), and Part E(c),(d).
\end{statement}

\begin{proof}
\emph{The exact descent map.}
The polynomial $X^{3^d}-17$ is Eisenstein at seventeen. Hence
$[F_d:\mathbb Q]=3^d$, $[M_d:F_d]=2$, and
$M_{d+1}/M_d$ is cyclic of degree three, generated by
$\sigma(\alpha_{d+1})=\zeta\alpha_{d+1}$.
The identity $c\sigma c=\sigma^{-1}$ holds on both generators.
Let $z\in F_d^\times$ and let $b\in M_{d+1}^\times$ satisfy
$N_{M_{d+1}/M_d}b=z$.
Then
\[
 w=\frac{z}{b\,c(b)}\in F_{d+1}^\times,
 \qquad N_{F_{d+1}/F_d}w=\frac{z^3}{z^2}=z.
 \tag{E65.14}
\]
Conversely, an actual pure norm remains a norm after adjoining
$\zeta$. This proves the equivalence of the two norm tests by an
explicit preimage map.

The same formulas hold in the multiplication completions.
For clarity, the actual field inclusion gives
\[
 (\widehat{M_d^\times}_3)^c=
        \widehat{F_d^\times}_3,
 \qquad
 v\longmapsto N_{M_d/F_d}(v)^{1/2}
 \tag{E65.15}
\]
as inverse on the fixed subgroup.
Indeed the inclusion followed by norm is squaring, which is an
automorphism of a pro-$3$ multiplication completion; for a fixed
$v$, inclusion of its norm is $v^2$.
This proves injectivity, surjectivity onto the stated fixed subgroup,
and continuity of both maps. The corresponding assertion for norm
quotients also holds before completion. Given a $c$-fixed class
$[a]$ in the cyclic norm quotient, whose exponent divides three,
the actual element $(a\,c(a))^2\in F_d^\times$ has class
$4[a]=[a]$. Injectivity follows from (E65.14).

\emph{Conjugation on the local norm coordinates.}
Identify the cyclic norm quotient with the relative Brauer group by
the crossed product $a\mapsto(M_{d+1}/M_d,\sigma,a)$.
The actual conjugation map sends this algebra to
$(M_{d+1}/M_d,\sigma^{-1},c(a))$, which is the inverse Brauer class
of $(M_{d+1}/M_d,\sigma,c(a))$.
Thus, when the generator $\sigma$ is kept fixed, $c$ acts on local
norm coordinates by negative permutation of the places.
The cyclic crossed-product identification, the exact relative
Brauer sequence and the Hasse norm theorem are the opened results
\cite[IV, Theorem3.14; VIII, \S4, the exact sequence preceding
Lemma4.1 and its proof; VIII, Theorem3.1]{E_MilneCFT}.
They identify the global cyclic norm quotient with the direct sum
of its nontrivial local norm groups, subject to the one sum-of-invariants
relation.

All these nontrivial groups have order three. On a place fixed by
$c$, a $c$-fixed tuple satisfies $a=-a$, hence has value zero.
On a pair of conjugate places its values are $a,-a$; their sum is
zero. It follows that the fixed quotient is exactly one independent
$\mathbb F_3$ for every conjugate pair at which the cyclic local
extension is nontrivial. Every finite assignment to those pairs is
realized: its sum of invariants is zero, the relative Brauer sequence
supplies a global cyclic norm class, and the averaging formula after
(E65.15) supplies an actual representative in $F_d^\times$.
This proves both the global sufficiency and the cokernel surjectivity
that will be used in (E65.4).

\emph{Identifying every relevant local pair.}
Away from three and seventeen, the relative polynomial
$X^3-\alpha_d$ is etale. If $q_v\equiv2\pmod3$, adjoining $\zeta$
gives the unramified quadratic completion, fixed by $c$.
Cubing is a bijection on $\mathbb F_{q_v}^\times$, so
$\alpha_d$ already has a cube root there. The pure local algebra is
the product of the base field and an unramified quadratic field.
Its norm is onto because of the first factor.
If $q_v\equiv1\pmod3$, the quadratic algebra splits into two
conjugate copies. On either copy, (E65.3) is the Kummer splitting
test for the unramified cubic extension. A nonzero symbol gives an
unramified cubic field, whose unit norm is onto and whose uniformizer
norm has valuation three. A zero symbol gives three copies of the
base field. This yields exactly the coordinates
$\gamma_{d,v}v(z)$ in (E65.4).

At seventeen, $X^{3^d}-17$ is Eisenstein both over $\mathbb Q_{17}$
and over its unramified quadratic extension $K_{17}$.
There is one prime in $M_d$ over seventeen, fixed by $c$.
Directly in the pure fields, the relative extension is tamely and
totally ramified of degree three, its residue field is
$\mathbb F_{17}$, and
$N(\alpha_{d+1})=\alpha_d$.
Cubing is a bijection on $\mathbb F_{17}^\times$ and on the
principal-unit pro-seventeen group. Hence all local valuations and
units are norms; this gives the same absence of a seventeen
coordinate by an explicit local map.
Infinite places contribute no $3$-primary norm quotient.

For three, Hensel's lemma applied at $23$ gives a root $\rho$ of
$X^3-17$ with $\rho\equiv23\pmod{27}$:
$v_3(23^3-17)=5>2v_3(3\cdot23^2)=2$.
There is only one such root in $\mathbb Q_3$, because this field has
no primitive third root of unity.
For $d\ge1$ one has the exact factorization
\[
 X^{3^d}-17=(X^m-\rho)
                   (X^{2m}+\rho X^m+\rho^2),\quad m=3^{d-1}.
 \tag{E65.16}
\]
After $X=Z-1$, both factors are Eisenstein at three.
Their reductions are $Z^m$ and $Z^{2m}$; their constant terms
are $-1-\rho$ and $1-\rho+\rho^2$, each of valuation one.
Consequently $F_d$ has precisely two primes over three, with
ramification indices $m$ and $2m$, both with residue field
$\mathbb F_3$.
The first completion has odd degree and therefore does not contain
$K_3$. The second contains it, since
$\alpha_d^m/\rho$ is a primitive third root of unity.
After adjoining $\zeta$, the first prime remains a single place
fixed by $c$, while the second becomes a conjugate pair.
Choose one of the latter pair by
$\alpha_d^m=\rho\zeta$.
At the next level the two degrees in (E65.16) are multiplied by
three. On the chosen pair, adjoining $\alpha_{d+1}$ is the cyclic
totally ramified cubic extension given by (E65.7).
Thus there is exactly one nontrivial conjugate pair over three for
every $d\ge1$.
At $d=0$, the local algebra after adjoining $\zeta$ splits completely,
because $17=\rho^3$ in $\mathbb Q_3$; there is no pair contributing
a norm condition.
This proves the list of all coordinates in (E65.4).

\emph{A finite presentation of the wild local norm group.}
The Eisenstein polynomial $Q_d$ proves that $\pi$ is a uniformizer,
$v_3(3)=e$, and the integer ring is $\mathbb Z_3[\pi]$.
For example, the last assertion follows directly by writing an
element in the basis $1,\pi,\ldots,\pi^{e-1}$: the valuations of
its nonzero terms are distinct modulo $e$, so integrality forces
all coefficients to be $3$-adic integers.
Since $3e\ge2e+1=n$, reduction modulo $\pi^n$ kills $27$.
It therefore depends only on $\rho\bmod27$ and is exactly (E65.6).
Each successive unit quotient
$(1+\pi^i\mathcal O)/(1+\pi^{i+1}\mathcal O)$ is the additive
$\mathbb F_3$ generated by $1+\pi^i$.
Successively removing the leading term gives the unique ordered
unit coordinates used in (E65.8).

The group of principal units modulo $\pi^n$ has a presentation with
commuting generators $g_i$ and the relations expressing $g_i^3$ in
the generators of higher index. These relations reduce every word
to exponents $0,1,2$, so the presented group has at most
$3^{n-1}$ elements; the actual unit filtration has exactly that many.
The presentation is therefore exact. An $\mathbb F_3$-valued
character of the unit group is precisely a vector annihilating the
first set of rows in (E65.8).

In the actual next completion, $\Pi=\alpha_{d+1}+1$ is a uniformizer
and its minimal polynomial over the base is the polynomial in
(E65.7). It is Eisenstein at $\pi$ and has norm $\pi$.
The target quotient of its integer ring by $\pi^n$ is exactly
$R'_d$, and $\Pi^{3n}=0$ there.
Its units are generated by $-1$ and
$1+\Pi^j$ for $1\le j<3n$.
Norms on this finite free ring are the displayed determinants.
Units congruent to one modulo $\pi^n$ in the source have norm
congruent to one modulo $\pi^n$ in the target, so these rows compute
the whole unit norm image in the finite quotient.

Finally $1+\pi^n\mathcal O$ is contained in the cube subgroup of
the base field: for $u$ in this group, Hensel's lemma applied to
$X^3-u$ at $1$ uses
$v(1-u)\ge2e+1>2v(3)$.
A cube in the base is the norm of its base-field cube root, viewed
in the cubic extension. Thus this deep unit group is already in
the actual norm image.
Local reciprocity gives index three for the norm of a cyclic cubic
local extension, and $N\Pi=\pi$ removes the valuation quotient
\cite[III, Theorem3.4 and its norm-index proof]{E_MilneCFT}.
The unit quotient computed by (E65.8) therefore has exactly order
three. This proves that its character nullspace is one-dimensional,
and proves necessity and sufficiency of the explicit character
(E65.9) at the wild prime.

\emph{Completion and all norm fibers.}
Let $B=F_{d+1}^\times$, let $I=N B\subset F_d^\times$, and let
$C=\ker(B\to I)$. In additive notation,
$0\to C\to B\to I\to0$ has $3$-torsion-free quotient $I$,
because $F_d$ is real and has only the roots of unity $\{1,-1\}$.
Thus $C\cap3^r B=3^r C$ for every $r$.
The exact finite-power quotient sequences have surjective transition
maps on the kernel quotients. Their inverse limit is exact and
identifies $\widehat C_3$ with the kernel of
$\widehat B_3\to\widehat I_3$.
Separately, $3F_d^\times\subset I$, because the norm of an included
base element is its cube. Hence
\[
 I\cap3^{r+1}F_d^\times\subset3^r I
                  \subset I\cap3^rF_d^\times.
 \tag{E65.17}
\]
These cofinal filtrations embed $\widehat I_3$ as the actual kernel
of the discrete exponent-three quotient in (E65.4).
This proves the completed norm test and the claimed kernel without
using compactness of the infinite multiplicative groups.

The group $C$ has no torsion, since the norm of $-1$ in an odd-degree
extension is $-1$.
The quotient $F_{d+1}^\times/\{1,-1\}$ is free abelian of countable
rank: its unit subgroup is finitely generated free, and its quotient
is the group of principal divisors, a subgroup of the free abelian
group on the finite primes. The latter is free, so the exact sequence
splits as abelian groups. The subgroup $C$ is therefore free abelian
of at most countable rank.
Its rank is infinite. Take infinitely many rational primes splitting
completely in the actual normal closure
$\mathbb Q(\mu_{3^{d+1}},\alpha_{d+1})/\mathbb Q$, using
\cite[VIII, Theorem7.4]{E_MilneCFT}.
This field contains $M_{d+1}=\mathbb Q(\zeta,\alpha_{d+1})$
under the specified root choices, so those primes split completely
in $F_{d+1}$ as well.
For each choose distinct primes $P,Q$ of $F_{d+1}$ above the same
prime of $F_d$. A class-number multiple of $P-Q$ is the divisor of
an element $y$. These elements have disjoint prime supports and are
independent. Their norms are units of $F_d$, a group of finite rank.
The kernel of the resulting map from a countably infinite free
abelian group to this finite-rank unit group has infinite rank and
lies in $C$.
Choose a free basis $u_i$ of $C$.
Completing $\bigoplus_i\mathbb Z u_i$ for multiplication by $3^r$
gives exactly (E65.10), with the displayed product formula.
This also specifies the topology of each nonempty fiber.

\emph{Exact finite certificates and the positional comparison.}
The retained local program constructs (E65.6) by integral Hermite
normal form and checks the index of every ideal $\pi^k$, $0\le k\le n$.
It computes all determinant norms in (E65.8), and independently
checks character additivity with the actual ring multiplication.
For $d=1,2,3$ the nullspace is one-dimensional as proved above;
the first two normalized vectors are
\[
 \mathbf w_1=(1,2,0,0),\qquad
 \mathbf w_2=(1,2,1,1,2,2,2,2,0,0,0,0).
 \tag{E65.18}
\]
The full degree-eighteen local vector, matrices and relations
are retained in the output. Substitution of
$\varepsilon=7\pi-25$ gives (E65.11).
The official PARI program separately certifies the two number fields
with \texttt{bnfcertify=1}, records their actual integer rings and
prime ideals, and computes the relative norms in (E65.12) through
the embedding $\alpha_1\mapsto\alpha_2^3$.
The Python verification recomputes the four norm determinants in
characteristic zero. Thus the unit basis, the sign row, and the
non-norm witness have independent exact certificates.
The actual prime contractions away from $3,17$ are also checked
against (E65.3), including the linear--quadratic local algebras when
$q_v\equiv2\pmod3$.

Finally E.62 factors the positional algebra as $F_d\times L_d$.
The root-cubing map acts on these factors by the actual relative
field norms in E.46. Hence (E65.13) is that same positional arrow,
with both components now evaluated.
Norm transitivity gives
\[
 N_{F_d/\mathbb Q}(N_d^{\rm p}u)\,
 N_{L_d/\mathbb Q}(N_d v)
   =N_{F_{d+1}/\mathbb Q}(u)\,
      N_{L_{d+1}/\mathbb Q}(v).
 \tag{E65.19}
\]
The congruence modulo cubes defining $J$ is therefore preserved
literally, as claimed.
\end{proof}
